Iterative self-training with reward curation, where a reward model selects the best candidate response at each generation for supervised fine-tuning, is widely expected to erode output diversity through compounding selection pressure. We test this expectation with a controlled factorial experiment spanning 7 configurations, 3 iterative generations, and both Low-Rank Adaptation (LoRA) and full fine-tuning. The central finding contradicts the prevailing concern: semantic diversity, measured by Vendi Score over sentence embeddings, remains within $\pm$1\% of its initial value across all curated conditions. However, this stability masks what we term \emph{metric-dependent collapse}. Format diversity degrades substantially, with opening unique ratio falling from 0.701 to 0.367 and near-duplicate start rates rising from 50\% to 84\%, while embedding-based metrics remain blind to this convergence. We further document a direction reversal: curation suppresses diversity under LoRA ($-3.40$ percentage points relative to a random baseline) but promotes it under full fine-tuning ($+0.63$ pp), a pattern that tracks parameter capacity rather than selection intensity. These results establish that reward curation does not cause semantic collapse in iterative self-training, that diversity audits relying solely on embedding-based measures miss significant structural homogenization, and that conclusions about curation effects depend critically on the training regime under which they are measured.
